\documentclass[10pt,a4paper]{amsart}
\usepackage[margin=19mm]{geometry}
\usepackage{amsmath,amssymb,mathtools}
\usepackage[scr=rsfs,cal=euler]{mathalfa}
\usepackage{xfrac}
\usepackage[unicode,hidelinks,bookmarks=false]{hyperref}
\newcommand{\ie}{i.e.\ }
\newcommand{\cf}{cf.\ }
\newcommand{\resp}{respectively\ }
\newcommand{\Subsection}{\subsection}
\providecommand{\nobreakdash}{\nobreak\hbox{-}}
\setlength{\emergencystretch}{2em}
\multlinegap0pt
\usepackage{amssymb}
\usepackage[shortlabels]{enumitem}
\usepackage{tikz}
\usepackage{xcolor}
\newtheorem{lemma}{Lemma}
\newtheorem{proposition}{Proposition}
\newcommand{\N}{\mathbb{N}}
\newcommand{\sgrboth}[2]{\mathsf{b}_{#1}(#2)}
\newcommand{\sgrleft}[2]{\mathsf{l}_{#1}(#2)}
\newcommand{\sgrright}[2]{\mathsf{r}_{#1}(#2)}
\title{Stochastic Sandpiles with Uniform Toppling Rule on the Line}
\author{David Beck-Tiefenbach, Robin Kaiser}
\date{Source: arXiv:2603.16304v1 (CC BY 4.0)}
\begin{document}
\maketitle

\noindent\emph{Source and licence.} Stochastic Sandpiles with Uniform Toppling Rule on the Line. Version arXiv:2603.16304v1 (CC BY 4.0), distributed by arXiv under the Creative Commons Attribution 4.0 International licence, http://creativecommons.org/licenses/by/4.0/, as recorded in the arXiv licence metadata for this exact version and shown on the abstract page https://arxiv.org/abs/2603.16304v1. Full licence notice: \texttt{licenses/arxiv-2603.16304-CC-BY-4.0.txt}.\par\smallskip

\noindent\textit{Editorial scope.} Counted unit: the article's proposition prop:hole-uniformity (source0.tex lines 381--386). The article's complete original proof of it is delivered with it (source0.tex lines 388--429, one \texttt{proof} environment of 42 lines). No mathematics is written, repaired or re-derived: the statement and the proof are the article's own lines, in the article's own notation, and no retained line is substituted or re-flowed.  Retained uncounted as the source-local context the proof needs: lines 349--363.  Nothing was omitted from any retained span, so the package carries no omission record.  No retained line contains a citation, so the wrapper carries no bibliography and no bibliography entry.  No retained line contains a figure environment, an image-inclusion command or drawing code, so no image is delivered and the package declares no asset dependency.  Editorial formatting only, in addition: an AMS \texttt{amsart} preamble that restates the article's own declarations and macros used by the retained lines, namely the environments \texttt{lemma}, \texttt{proposition} on separate counters, together with the standard packages those lines need.  The retained environments print in this document as Lemma 1, Proposition 1 in order of appearance, whereas the article prints the same items with the numbering of its own full document.  The complete native source of the article as distributed in the arXiv e-print for this exact version is bundled unchanged as \texttt{sources/196533/source0.tex}.
\begin{lemma}[Hole Recurrence Relations]\label{lemma:hole-recurrence}
    For $n \ge 2$, the hole probabilities satisfy the following relations. 
    For $i < n$, we have the bulk recurrence:
    \begin{equation*}
        \mathsf{h}_n^j(i) = \sgrright{n-1}{i}\mathsf{h}_n^j(n) + \sum_{x=1}^{j-1} \mathsf{h}_{n-1}^x(i)\mathsf{h}_{n-x}^{j-x}(n-x) + \mathsf{h}_{n-1}^j(i)\sgrright{n-j}{n-j}.
    \end{equation*}
    For $i=n$, we have the boundary recurrence:
    \begin{equation*}
        \mathsf{h}_n^j(n) = 
        \begin{cases}
            \frac{1}{3}\mathsf{h}_n^j(n-1) + \frac{1}{3}\mathsf{h}_{n-1}^j(n-1) & \text{if } j < n, \\
            \frac{1}{3}\mathsf{h}_n^n(n-1) + \frac{1}{3}\sgrleft{n-1}{n-1} & \text{if } j = n.
        \end{cases}
    \end{equation*}
\end{lemma}

\begin{proposition}[Uniformity of the Hole]\label{prop:hole-uniformity}
    For all $n \in \N$ and $i, j \in C_n$, the probability of a hole forming at $j$ is independent of $j$ and is given by:
    \begin{equation*}
        \mathsf{h}_n^j(i) = \frac{\sgrboth{n}{i}}{n} = \frac{2i(n-i+1)}{n(n+1)(n+2)}.
    \end{equation*}
\end{proposition}

\begin{proof}
    We prove this by showing that $g_n^j(i) := \frac{\sgrboth{n}{i}}{n}$ uniquely satisfies the recurrence relations of Lemma \ref{lemma:hole-recurrence}.
    
     We first show that the boundary and bulk recurrence relations uniquely determine the entire sequence of probabilities. We prove this by induction on $n$. The base case $n=1$ is uniquely determined by the toppling rules. Assume the sequence is uniquely determined for all system sizes $m < n$. For system size $n$, evaluating the bulk recurrence at $n-1$ yields $$\mathsf{h}_n^j(n-1) = \sgrright{n-1}{n-1}\mathsf{h}_n^j(n) + K_j,$$ where 
     $$K_j:=\sum_{x=1}^{j-1} \mathsf{h}_{n-1}^x(n-1)\mathsf{h}_{n-x}^{j-x}(n-x) + \mathsf{h}_{n-1}^j(n-1)\sgrright{n-j}{n-j},$$
     depends only on systems of size $m < n$ and the location of the hole $j$, and is therefore fixed. Substituting this into the boundary recurrence yields a linear equation of the form
    $$ \mathsf{h}_n^j(n) = \frac{1}{3} \Big( \sgrright{n-1}{n-1}\mathsf{h}_n^j(n) + K_j \Big) + C_j, $$
    where
    $$j<n:C_j:=\frac{1}{3}\mathsf{h}_{n-1}^j(n-1),\hspace{1cm}C_n:=\frac{1}{3}\sgrleft{n-1}{n-1},$$
    is another fixed constant independent of the hole probabilities of system size $n$. Because the coefficient $\sgrright{n-1}{n-1}/3$ is strictly less than $1$, this equation has a strictly unique scalar solution for $\mathsf{h}_n^j(n)$. With the boundary value uniquely fixed, the bulk recurrence explicitly and uniquely determines the remaining values $\mathsf{h}_n^j(i)$ for $i < n$. Thus, if $g_n^j(i)$ satisfies the recurrences, then $g_n^j(i)= \mathsf{h}_n^j(i)$ for all $n\in\N$ and all $i,j\in C_n$.
    
    We will now verify that our proposed function $(g_n^j(i))_{i,j\in C_n}$ satisfies the boundary recurrence. The left-hand side of the equation simplifies to
    $$g_n^j(n) = \frac{\sgrboth{n}{n}}{n} = \frac{2n}{n(n+1)(n+2)} = \frac{2}{(n+1)(n+2)}.$$
    Now, calculating the right-hand side for $j < n$ gives
    \begin{align*}
        \frac{1}{3} \bigg( g_n^j(n-1) + g_{n-1}^j(n-1)\bigg) &= \frac{1}{3} \bigg(\frac{4(n-1)}{n(n+1)(n+2)} + \frac{2(n-1)}{(n-1)n(n+1)}\bigg)\\ &= \frac{1}{3} \frac{4(n-1) + 2(n+2)}{n(n+1)(n+2)}  = \frac{1}{3} \frac{6n}{n(n+1)(n+2)}  = \frac{2}{(n+1)(n+2)}. 
    \end{align*} 
    Similarly, the right-hand side for $j = n$ gives
    \begin{align*}
        \frac{1}{3} \bigg( g_n^n(n-1) + \sgrleft{n-1}{n-1}\bigg) &= \frac{1}{3} \bigg(\frac{4(n-1)}{n(n+1)(n+2)} + \frac{2}{n(n+1)}\bigg)\\ &= \frac{2}{(n+1)(n+2)}. 
    \end{align*} 
    This shows that $(g_n^j(i))_{i,j\in C_n}$ satisfies the boundary recurrence.
    
    Next, we will verify that our proposed function also satisfies the bulk recurrence. We substitute $g$ into the right-hand side for $i < n$. Notice that by definition, $g_{n-1}^x(i)$ does not depend on $x$. Thus, we can factor it out of the sum $S_j := \sum_{x=1}^{j-1} g_{n-1}^x(i)g_{n-x}^{j-x}(n-x)$ to obtain
    \begin{align*}
        S_j &= \frac{\sgrboth{n-1}{i}}{n-1} \sum_{x=1}^{j-1} \frac{2(n-x)}{(n-x)(n-x+1)(n-x+2)} = \frac{\sgrboth{n-1}{i}}{n-1} \sum_{x=1}^{j-1} \left( \frac{2}{n-x+1} - \frac{2}{n-x+2} \right) \\&= \frac{\sgrboth{n-1}{i}}{n-1} \left( \frac{2}{n-j+2} - \frac{2}{n+1} \right),
    \end{align*}
    Combining $S_j$ with $g_{n-1}^j(i)\sgrright{n-j}{n-j}$, gives
    \begin{align*}
        S_j + g_{n-1}^j(i)\sgrright{n-j}{n-j} &= \frac{\sgrboth{n-1}{i}}{n-1} \left[ \left( \frac{2}{n-j+2} - \frac{2}{n+1} \right) + \frac{n-j}{n-j+2} \right] \\
        &= \frac{\sgrboth{n-1}{i}}{n-1} \left[ \frac{n-j+2}{n-j+2} - \frac{2}{n+1} \right] = \frac{\sgrboth{n-1}{i}}{n-1} \frac{n-1}{n+1} = \frac{2i(n-i)}{n(n+1)^2}.
    \end{align*}
    Finally, adding $\sgrright{n-1}{i}g_n^j(n)$ we see that the right-hand side of the bulk recurrence simplifies to
    \begin{align*}
        \frac{i(i+1)}{n(n+1)}   \frac{2}{(n+1)(n+2)}  &+ \frac{2i(n-i)}{n(n+1)^2}
        = \frac{2i}{n(n+1)^2(n+2)} \Big[ i + 1 + (n-i)(n+2) \Big] \\
        &= \frac{2i}{n(n+1)^2(n+2)} \Big[ n^2 + 2n + 1 - i(n+1) \Big] \\
        &= \frac{2i(n+1)(n-i+1)}{n(n+1)^2(n+2)} = \frac{2i(n-i+1)}{n(n+1)(n+2)} = g_n^j(i),
    \end{align*}
    as desired.
    Since $(g_n^j(i))_{i,j\in C_n}$ satisfies both recurrences, and the solution to the recurrences is unique, we obtain $\mathsf{h}_n^j(i) = g_n^j(i)$ for all $n\in\N$ and all $i,j\in C_n$, concluding the proof.
\end{proof}

\end{document}
